
\subsubsection{Key Findings}

\textbf{Finding 1: Truthfulness is multi-dimensional.} TruthfulQA, HaluEval, and FactScore measure partially independent dimensions. No single benchmark captures the full picture. Evaluating a model on one benchmark alone provides incomplete information about its reliability.

\textbf{Finding 2: Knowledge does not guarantee misconception resistance.} The divergent profile analysis (4 models with high MMLU but low TruthfulQA) demonstrates that strong factual knowledge does not guarantee avoiding misconceptions. Models can learn facts and falsehoods from training data simultaneously.

\textbf{Finding 3: FactScore captures an orthogonal dimension.} The near-zero correlations between FactScore and other benchmarks (r $\approx$ 0) suggest that atomic factual precision in long-form generation involves capabilities largely independent of misconception detection or hallucination classification.

\subsubsection{Theoretical Interpretation}

The findings support a mechanistic view of truthfulness comprising three stages:

\begin{enumerate}
\item \textbf{Knowledge Selection:} Choosing to output factual rather than popular misconceptions (TruthfulQA)
\item \textbf{Coherence Maintenance:} Maintaining consistency and avoiding invented details during generation (HaluEval)
\item \textbf{Fine-Grained Precision:} Ensuring atomic claims are verifiable against external sources (FactScore)
\end{enumerate}

A model can fail at any stage independently.

\subsubsection{Unexpected Finding: Near-Zero FactScore Correlations}

The near-zero or negative correlations between FactScore and other benchmarks warrant consideration. Two explanations:

\begin{enumerate}
\item \textbf{Task format difference:} FactScore uses long-form biography generation with atomic decomposition and retrieval verification---fundamentally different from MC questions or binary classification.
\end{enumerate}

\begin{enumerate}
\item \textbf{Genuinely orthogonal construct:} Atomic factual precision may represent a distinct capability from holistic response quality.
\end{enumerate}

These explanations are not mutually exclusive.

\subsubsection{Limitations}

\textbf{Model population scope.} Analysis includes only open-source, decoder-only models at 7B--70B scale with English evaluation. Correlation structure may differ for proprietary models, encoder-decoder architectures, sub-7B or 100B+ scales, or multilingual evaluation. Open-source models dominate research; findings provide an actionable baseline.

\textbf{FactScore proxy methodology.} Due to computational constraints, proxy FactScore estimates were used rather than full atomic decomposition on all 50 models. Proxy methodology was validated against available official scores; the key finding (near-zero correlation) is robust across estimation approaches.

\textbf{Intra-HaluEval correlation below threshold.} HaluEval internal correlation (r=0.65) fell below predicted r > 0.7. The relative pattern holds (intra > inter: 0.65 > 0.16). This finding itself indicates that even within HaluEval, subtasks target somewhat different coherence failure modes.

\subsubsection{Implications for Practice}

\textbf{Evaluation guidance:} Model selection should assess all three dimensions. Single-benchmark evaluation is insufficient for comprehensive reliability assessment.

\textbf{Intervention design:} Improving one dimension may not transfer to others. Targeted interventions may be needed for comprehensive improvement.
